% Part: many-valued-logic
% Chapter: three-valued-logics
% Section: multiple-designation

\documentclass[../../../include/open-logic-section]{subfiles}

\begin{document}

\olfileid[hi]{mvl}{thr}{mul}

\olsection{केवल $\True$ को निर्दिष्ट न करना}

अब तक हमने जितने तर्कशास्त्र देखे हैं, उन सभी में निर्दिष्ट सत्य-मूल्यों का
समुच्चय $V^+ = \{\True\}$ था; अर्थात् कोई चीज तभी सत्य मानी जाती है जब
उसका सत्य-मूल्य $\True$ हो। किंतु किसी चीज को केवल $\False$ न होने पर भी
सत्य माना जा सकता है। तब, उदाहरणार्थ, आव्यूह में $V^+ = \{\True,
\Undef\}$ निर्धारित करके एक तर्कशास्त्र प्राप्त होगा।

\begin{defn}
\emph{विरोधाभास का तर्कशास्त्र}~$\LogLP$ निम्नलिखित आव्यूह से परिभाषित होता
है:
\begin{enumerate}
  \item $\lnot$, $\land$, $\lor$, $\lif$ वाली मानक प्रतिज्ञप्तिक भाषा
  $\Lang L_0$।
  \item सत्य-मूल्यों का समुच्चय $V = \{\True, \Undef, \False\}$।
  \item $\True$ और $\Undef$ निर्दिष्ट हैं, अर्थात्
  $V^+ = \{\True, \Undef\}$।
  \item सत्य-फलन प्रबल क्लिनी तर्कशास्त्र के समान हैं।
\end{enumerate}
\end{defn}

\begin{defn}
हॉल्डेन का \emph{निरर्थकता का तर्कशास्त्र}~$\LogHal$ निम्नलिखित आव्यूह से
परिभाषित होता है:
\begin{enumerate}
  \item $\lnot$, $\land$, $\lor$, $\lif$ और एक $1$-स्थानिक संयोजक~$+$
  वाली मानक प्रतिज्ञप्तिक भाषा $\Lang L_0$।
  \item सत्य-मूल्यों का समुच्चय $V = \{\True, \Undef, \False\}$।
  \item $\True$ और $\Undef$ निर्दिष्ट हैं, अर्थात्
  $V^+ = \{\True, \Undef\}$।
  \item सत्य-फलन दुर्बल क्लिनी तर्कशास्त्र के समान हैं, और इनके साथ
  ``निरर्थक है'' संकारक भी है:
  \begin{center}
    \begin{tabular}{c|c} 
    $\tf{+}$ & \\ 
    \hline  
    $\True$ & $\False$ \\ 
    $\Undef$ & $\True$ \\
    $\False$ & $\False$ 
    \end{tabular}
  \end{center}
\end{enumerate}
\end{defn}

जिन क्लिनी तर्कशास्त्रों के साथ इनकी सत्य-सारणियाँ समान हैं, उनके विपरीत
इनमें सर्वसत्यताएँ \emph{होती हैं}।

\begin{prop}\ollabel{prop:LP-taut-CL} $\LogLP$ की सर्वसत्यताएँ वही हैं जो
  शास्त्रीय प्रतिज्ञप्तिक तर्कशास्त्र की सर्वसत्यताएँ हैं।
\end{prop}

\begin{proof}
  \olref[syn][sub]{prop:mvl-cl} के अनुसार, यदि $\Entails[\LogLP] !A$ तो
  $\Entails[\LogCL] !A$। विलोम दिखाने के लिए हम सिद्ध करते हैं कि यदि कोई
  !!{valuation}~$\pAssign v\colon \PVar \to \{\False, \True, \Undef\}$
  ऐसा है कि $\pValue v(!A)[\LogKs] = \False$, तो कोई
  !!{valuation}~$\pAssign {v'}\colon \PVar \to \{\False, \True\}$ ऐसा है
  कि $\pValue {v'}(!A)[\LogCL] = \False$। इससे $\LogLP$ के लिए परिणाम
  स्थापित हो जाता है, क्योंकि $\LogKs$ और $\LogLP$ के अभिलाक्षणिक सत्य-फलन
  समान हैं और $\False$, $\LogLP$ का एकमात्र अनिर्दिष्ट सत्य-मूल्य है
  ($\LogLP$ और $\LogKs$ में केवल यही अंतर है)। अतः यदि
  $\Entails/[\LogLP] !A$, तो किसी !!{valuation}~$\pAssign v$ के लिए
  $\pValue v(!A)[\LogLP] = \pValue v[\LogKs](!A) = \False$। हमारे दावे के
  अनुसार $\pValue{v'}[\LogCL](!A) = \False$, अर्थात्
  $\Entails/[\LogCL] !A$।
  
  दावा स्थापित करने के लिए पहले $\pAssign {v'}$ को इस प्रकार परिभाषित करते
  हैं:
  \[
    \pAssign {v'}(p) = 
    \begin{cases}
    \True & \text{यदि } \pAssign {v}(p) \in \{\True, \Undef\}\\
    \False & \text{अन्यथा}
  \end{cases}
  \]
  अब हम $!A$ पर आगमन से दिखाते हैं कि (a)~यदि
  $\pValue v(!A)[\LogKs] = \False$, तो
  $\pValue {v'}(!A)[\LogCL] = \False$, और (b)~यदि
  $\pValue v(!A)[\LogKs] = \True$, तो
  $\pValue {v'}(!A)[\LogCL] = \True$।
  \begin{enumerate}
    \item आगमन-आधार: $!A \ident p$।
    \olref[syn][val]{defn:pValue} के अनुसार,
    $\pValue v(!A)[\LogKs] = \pAssign v(p) =
    \pValue {v'}(!A)[\LogCL]$, जिससे (a) और (b) दोनों निष्कर्ष निकलते हैं।
    
    आगमन-चरण के लिए निम्नलिखित स्थितियों पर विचार करें:
    \item $!A \ident \lnot !B$।
    \begin{enumerate}
      \item मान लें कि $\pValue v(\lnot!B)[\LogKs] = \False$।
      $\tf{\lnot}[\LogKs]$ की परिभाषा से
      $\pValue v(!B)[\LogKs] = \True$। आगमन-परिकल्पना के (b) मामले से
      $\pValue {v'}(!B)[\LogCL] = \True$ मिलता है, अतः
      $\pValue {v'}(\lnot !B)[\LogCL] = \False$।
      \item मान लें कि $\pValue v(\lnot!B)[\LogKs] = \True$।
      $\tf{\lnot}[\LogKs]$ की परिभाषा से
      $\pValue v(!B)[\LogKs] = \False$। आगमन-परिकल्पना के (a) मामले से
      $\pValue {v'}(!B)[\LogCL] = \False$ मिलता है, अतः
      $\pValue {v'}(\lnot !B)[\LogCL] = \True$।
    \end{enumerate}

    \item $!A \ident (!B \land !C)$।
    \begin{enumerate}
      \item मान लें कि $\pValue v(!B \land !C)[\LogKs] = \False$।
      $\tf{\land}[\LogKs]$ की परिभाषा से
      $\pValue v(!B)[\LogKs] = \False$ या
      $\pValue v(!B)[\LogKs] = \False$। आगमन-परिकल्पना के (a) मामले से
      $\pValue {v'}(!B)[\LogCL] = \False$ या
      $\pValue {v'}(!C)[\LogCL] = \False$ मिलता है, अतः
      $\pValue {v'}(!B \land !C)[\LogCL] = \False$।
      \item मान लें कि $\pValue v(!B \land !C)[\LogKs] = \True$।
      $\tf{\land}[\LogKs]$ की परिभाषा से
      $\pValue v(!B)[\LogKs] = \True$ और
      $\pValue v(!B)[\LogKs] = \True$। आगमन-परिकल्पना के (b) मामले से
      $\pValue {v'}(!B)[\LogCL] = \True$ और
      $\pValue {v'}(!C)[\LogCL] = \True$ मिलता है, अतः
      $\pValue {v'}(!B \land !C)[\LogCL] = \True$।
    \end{enumerate}
  \end{enumerate}

  अन्य दो मामले इसी प्रकार हैं और अभ्यास के लिए छोड़े गए हैं। वैकल्पिक रूप
  से, ऊपर का प्रमाण केवल $\lnot$ और $\land$ वाली सभी !!{formula}ों के लिए
  परिणाम स्थापित करता है। अब इस तथ्य का उपयोग किया जा सकता है कि $\LogKs$
  और $\LogCL$ दोनों में, प्रत्येक $\pAssign v$ के लिए,
  $\pValue v(!B \lor !C) = \pValue v(\lnot(\lnot!B \land \lnot !C))$
  और $\pValue v(!B \lif !C) =
  \pValue v(\lnot(!B \land \lnot !C))$।
\end{proof}

\begin{prob}
  \olref[mvl][thr][mul]{prop:LP-taut-CL} का प्रमाण पूरा कीजिए; अर्थात् उन
  मामलों में (a) और (b) स्थापित कीजिए जिनमें
  $!A \ident (!B \lor !C)$ और $!A \ident (!B \lif !C)$।
\end{prob}

\begin{prob}
सिद्ध कीजिए कि प्रत्येक शास्त्रीय सर्वसत्यता $\LogHal$ में सर्वसत्यता है।
\end{prob}

यद्यपि इनकी सर्वसत्यताएँ शास्त्रीय तर्कशास्त्र जैसी ही हैं, इनके परिणाम
संबंध भिन्न हैं। उदाहरणार्थ, $\LogLP$ इस अर्थ में \emph{परासंगत} है कि
$\lnot p, p \Entails/ q$; अतः विस्फोट सिद्धांत
$\lnot !A, !A \Entails !B$ सामान्यतः मान्य नहीं है। (यह $!A$ और $!B$ के
कुछ मामलों में मान्य है, जैसे जब $!B$ सर्वसत्यता हो।)

\begin{prob}
  निम्नलिखित में से कौन-से संबंध (a)~$\LogLP$ में और (b)~$\LogHal$ में
  मान्य हैं? प्रत्येक के लिए सत्य-सारणी दीजिए।
  \begin{enumerate}
    \item $p, p \lif q \Entails q$
    \item $\lnot q, p \lif q \Entails \lnot p$
    \item $p \lor q, \lnot p \Entails q$
    \item $\lnot p, p \Entails q$
    \item $p \Entails p \lor q$
    \item $p \lif q, q\lif r \Entails p \lif r$
  \end{enumerate}
\end{prob}

यदि $\LogLuk[3]$ में $\Undef$ को निर्दिष्ट कर दिया जाए तो क्या होगा?

\begin{defn}
  \emph{त्रिमूल्य आर-मिंगल}~$\LogRM[3]$ निम्नलिखित आव्यूह से परिभाषित होता
  है:
  \begin{enumerate}
    \item $\lfalse$, $\lnot$, $\land$, $\lor$, $\lif$ वाली मानक
    प्रतिज्ञप्तिक भाषा $\Lang L_0$।
    \item सत्य-मूल्यों का समुच्चय $V = \{\True, \Undef, \False\}$।
    \item $\True$ और $\Undef$ निर्दिष्ट हैं, अर्थात्
    $V^+ = \{\True, \Undef\}$।
    \item सत्य-फलन लुकासिएविच तर्कशास्त्र~$\LogLuk[3]$ के समान हैं।
  \end{enumerate}
\end{defn}

\begin{prob}
  निम्नलिखित में से कौन-से संबंध $\LogRM[3]$ में मान्य हैं?
  \begin{enumerate}
    \item $p, p \lif q \Entails q$
    \item $p \lor q, \lnot p \Entails q$
    \item $\lnot p, p \Entails q$
    \item $p \Entails p \lor q$
  \end{enumerate}
\end{prob}

कभी-कभी केवल निर्दिष्ट मूल्यों को बदलने से भिन्न सत्य-सारणियाँ एक ही
तर्कशास्त्र (परिणाम संबंध) उत्पन्न कर सकती हैं। उदाहरणार्थ, गोडेल
तर्कशास्त्र में $V^+ = \{\True\}$ के स्थान पर
$V^+ = \{\True, \Undef\}$ लेने पर ऐसा होता है।

\begin{prop}\ollabel{prop:gl-udes}
  $V = \{\False, \Undef,\True\}$, $V^+=\{\True, \Undef\}$ और त्रिमूल्य
  गोडेल तर्कशास्त्र के सत्य-फलनों वाला आव्यूह शास्त्रीय तर्कशास्त्र को
  परिभाषित करता है।
\end{prop}

\begin{proof}
  अभ्यास।
\end{proof}

\begin{prob}
  \olref[mvl][thr][mul]{prop:gl-udes} को इस प्रकार सिद्ध कीजिए: गोडेल
  तर्कशास्त्र की तरह, पर $V^+=\{\True,\Undef\}$ के साथ परिभाषित
  तर्कशास्त्र~$\Log L$ के लिए यदि $\Gamma \Entails/[\Log L] !B$, तो
  $\Gamma \Entails/[\LogCL] !B$।
  \olref[mvl][thr][mul]{prop:LP-taut-CL} के विचारों का उपयोग कीजिए, किंतु
  गुण (a) और (b) सिद्ध करने के स्थान पर दिखाइए कि
  $\pValue v(!A)[\LogGod] = \False$ तभी जब
  $\pValue {v'}(!A)[\LogCL] = \False$ (और इसलिए
  $\pValue v(!A)[\LogGod] \in \{\True,\Undef\}$ तभी जब
  $\pValue {v'}(!A)[\LogCL] = \True$)। समझाइए कि इससे प्रतिज्ञप्ति क्यों
  स्थापित होती है।
\end{prob}

\end{document}
